% !TEX program = XeLaTeX
% glossaries 多类型测试（main + acronym + symbols + 自定义）
% 覆盖 detection.py 中 _NGLOSSARY_PATTERN 多类型发现路径
\documentclass{article}
\usepackage{ctex}
\usepackage{glossaries}

% 先创建术语表类型，再定义条目
\makeglossaries

% 术语表条目
\newglossaryentry{latex}{
    name={\LaTeX},
    description={一种基于 TeX 的高质量排版系统},
    sort={latex}
}
\newglossaryentry{tex}{
    name={\TeX},
    description={由 Donald Knuth 设计的数学排版系统},
    sort={tex}
}

% 首字母缩写条目（glossaries 内置 acronym 类型）
\newacronym{fsm}{FSM}{有限状态机}
\newacronym{cfg}{CFG}{上下文无关文法}
\newacronym{nlp}{NLP}{自然语言处理}

% 打印所有术语表（glossaries 会自动发现 main + acronym 两套）
\begin{document}

\section{多术语表测试}

本文档同时包含主术语表和缩写表，测试 glossaries 多类型发现能力。

\LaTeX \gls{latex} 和 \TeX \gls{tex} 是常用的排版系统。
在计算机科学中，有限状态机 \acrshort{fsm}、上下文无关文法 \acrfull{cfg}
和自然语言处理 \gls{nlp} 都是重要概念。

\printglossary
\printglossary[type=\acronymtype]

\end{document}
